\documentclass[a4paper, 11pt]{article}
\usepackage{geometry}
\usepackage{url}
\usepackage{color}


\linespread{1}

\geometry{a4paper,top=3cm,left=3cm,right=2.5cm,bottom=2cm}

\usepackage{hyperref}
\hypersetup{colorlinks,linkcolor=black,citecolor=blue}

% Additional packages required for tables
\usepackage{booktabs}
\usepackage{enumitem}


\title{\textbf{Laboratory assignment} \\[1ex] \large \textbf{Component 4:} Problem Definition and MAS Specification for SP2}

\author{\textbf{Authors:} Fusarau Ioan, Dreghici Bogdan\\ \textbf{Group:} ICA 1}

\date{April 28, 2026}

\begin{document}

\maketitle

\section{Problem Definition}

The objective of the second software project (SP2) is to design and
implement, \emph{from scratch}, a Multi-Agent System (MAS) that simulates
collaborative wildfire suppression on a discrete two-dimensional grid.
The system models a forest (or building floor) discretised as a grid of
cells, in which a small number of cells are initially set on fire. Fire
propagates stochastically over time to neighbouring vegetation, while a
configurable team of autonomous firefighter agents moves on the grid,
locating burning cells and extinguishing them.

The problem is interesting from a Distributed Artificial Intelligence
(DAI) perspective for three reasons. First, the global goal containing
the fire before it spreads out of control cannot be achieved
efficiently by a single agent: as the burning perimeter grows, only a
distributed team can react in time. Second, the agents must coordinate
implicitly, avoiding redundant work (e.g., two firefighters running to the
same cell) without a central commander. Third, the environment is
inherently dynamic and partially observable: the world changes between
agent decisions, and each agent perceives only a local neighbourhood.

A simulation run is considered \emph{successful} if all fires are
extinguished before a configurable damage threshold (e.g., 40\% of the
grid burnt) is reached. Otherwise, the run is considered a \emph{failure}.
The system reports per-tick statistics (number of burning cells, number of
extinguished cells, agent positions) and a final outcome summary.

\section{High-Level MAS Specification}

\subsection{Inputs and Outputs}

\paragraph{Inputs to the MAS.}
\begin{itemize}[leftmargin=*]
\item \textbf{Grid configuration:} grid dimensions \texttt{(rows, cols)}.
\item \textbf{Initial fire seeds:} a list of \texttt{(row, col)} cells set
on fire at $t=0$.
\item \textbf{Number of agents} $N$ (configurable, $N \geq 1$).
\item \textbf{Fire spread probability} $p \in [0, 1]$ --- the probability
that a burning cell ignites each unburnt orthogonal neighbour per tick.
\item \textbf{Maximum simulation length} (in ticks) and the
\textbf{damage threshold} for failure.
\item \textbf{Strategy selector} for the firefighter agents
(e.g., \emph{nearest-fire greedy}, \emph{region partitioning}).
\end{itemize}

\paragraph{Outputs of the MAS.}
\begin{itemize}[leftmargin=*]
\item \textbf{Per-tick simulation log:} grid state (each cell as
\textsc{intact}, \textsc{burning}, \textsc{extinguished}, or
\textsc{ash}); agent positions and current targets; claim table.
\item \textbf{Aggregate statistics:} total burnt area, number of
extinguishing actions per agent, average response time per fire.
\item \textbf{Final outcome:} \texttt{SUCCESS} (all fires out below the
damage threshold) or \texttt{FAILURE} (threshold exceeded or simulation
timeout).
\end{itemize}

\subsection{Types of Agents}

The MAS contains two types of entities:

\begin{enumerate}
\item \textbf{Firefighter Agent} (active, instantiated $N$ times):
an autonomous, mobile, reactive agent whose role is to perceive the local
grid, select a target burning cell, and extinguish it. Each firefighter
runs in its own thread, executing concurrently with the others.
\item \textbf{Environment Agent} (single instance): a passive coordinator
that maintains the authoritative grid state, advances fire propagation
each tick, and provides a \emph{blackboard} interface through which
firefighter agents read and write shared information (current grid,
claimed targets).
\end{enumerate}

\subsection{Specification of Agents}

\subsubsection{Firefighter Agent}

\begin{tabular}{p{3cm}p{10.5cm}}
\toprule
\textbf{Inputs}   & Local perception of the grid (cells within a configurable visibility radius), the current claim table, and the agent's own position. \\
\textbf{Outputs}  & Movement actions (\textsc{move-north}, \textsc{move-south}, \textsc{move-east}, \textsc{move-west}, \textsc{stay}); the \textsc{extinguish} action when standing on a burning cell; updates to the claim table. \\
\textbf{Task}     & Locate burning cells, claim a target on the blackboard, navigate to it, and extinguish it. Iterate until no burning cells remain or the simulation ends. \\
\bottomrule
\end{tabular}

\subsubsection{Environment Agent}

\begin{tabular}{p{3cm}p{10.5cm}}
\toprule
\textbf{Inputs}   & The current grid state; the action requests submitted by firefighter agents during the current tick. \\
\textbf{Outputs}  & Updated grid state at $t+1$: stochastic fire propagation applied to the neighbours of every burning cell, and the effects of each firefighter's action committed atomically. \\
\textbf{Task}     & Maintain a globally consistent world state, advance the simulation clock, enforce action validity, and detect terminal conditions (success / failure). \\
\bottomrule
\end{tabular}

\subsection{Communications among Agents}

Communication is realised through a \textbf{shared blackboard}, which is
the canonical communication model for spatial multi-agent simulations and
is one of the two models recommended in the project requirements (the
alternative being explicit message passing).

\begin{itemize}[leftmargin=*]
\item \textbf{Medium.} A central \texttt{Blackboard} object owned by the
Environment Agent contains: (i) the grid state, (ii) the agent
position table, and (iii) the \emph{claim table}, mapping each claimed
burning cell to the firefighter that committed to extinguish it.
\item \textbf{Synchronisation.} The blackboard is protected by a mutex
(reentrant lock). Reads return a consistent snapshot; writes are atomic.
This allows firefighter threads to coordinate without explicit
peer-to-peer messages.
\item \textbf{Coordination protocol.} Before moving toward a target,
a firefighter \emph{publishes a claim} on the blackboard. Other
firefighters see the claim during their own perception step and
exclude already-claimed cells from their candidate target set,
preventing redundant work. Claims expire when the target is either
extinguished or no longer burning.
\item \textbf{Implicit message types.} Although no FIPA-ACL message bus
is used, the interactions can be mapped onto speech acts for
documentation purposes: a claim corresponds to an \textsc{inform} act
(``I am committing to cell $X$''); a perception of a competitor's claim
corresponds to a passive \textsc{inform} reception; a successful
extinguishing corresponds to an \textsc{inform-done} act that frees the
claim.
\item \textbf{Configurability.} The communication model is encapsulated
behind a \texttt{CommunicationLayer} interface, so a future variant of
the system can replace the blackboard with explicit message passing
(e.g., per-agent inboxes implemented with thread-safe queues) without
touching the agent logic.
\end{itemize}

\subsection{Agents' Role within the MAS}

The firefighter agents collectively realise the system's primary goal:
\emph{containing and extinguishing the fire}. They are heterogeneous only
in identity (each has a unique ID) and homogeneous in behaviour, which
makes the team symmetric and easily scalable to any $N$. The Environment
Agent plays the supporting role of \emph{world simulator and arbiter};
it does not pursue an autonomous goal, but it ensures that the
firefighters' concurrent actions resolve into a coherent, conflict-free
shared reality. Together, the two agent types form a cooperative,
non-competitive MAS in which the only adversary is the dynamics of the
fire itself.

\section{Conclusion}

This document defines the problem to be solved by SP2 collaborative
distributed wildfire suppression and provides a high-level
specification of the corresponding multi-agent system. The specification
identifies the system's inputs and outputs, the two types of agents
involved (firefighter and environment), the precise input/output/task
contract for each, and the blackboard-based communication protocol that
enables decentralised coordination. The conceptual analysis (PAGES, PEAS,
environment properties), the architectural design, and the Python
implementation will be developed in the subsequent components.

\end{document}
